\SourceLeaf{444}%
folgt also der wichtige Lehrsatz, dass
\begin{quote}
»bei allen Flächen, die sich in eine Ebene abwickeln lassen, das Krümmungsmass überall verschwindet«
\end{quote}
\noindent oder dass
\[
\left(\frac{\partial\partial z}{\partial x.\partial y}\right)^{2}-\left(\frac{\partial\partial z}{\partial x^{2}}\right)\left(\frac{\partial\partial z}{\partial y^{2}}\right)=0
\]
wird, welches Criterium sonst aus andern Gründen abgeleitet wird, obwohl, wie uns scheint, nicht mit der zu wünschenden Evidenz. Es ist klar, dass bei allen solchen Flächen die Zenithe aller Punkte keinen Raum ausfüllen können, und daher alle in einer Linie liegen.

\begin{center}17.\end{center}
\noindent\hspace*{1.2em}Von einem gegebenen Punkte auf einer krummen Fläche wollen wir unzählige kürzeste Linien auf dieser Fläche ausgehen lassen, die von einander durch den Winkel unterschieden werden sollen, welchen ihr erstes Element mit dem ersten Element Einer \textls[100]{bestimmten} kürzesten Linie macht: dieser Winkel soll $\theta$ heissen. Es soll ferner $s$ die [von dem gegebenen Punkte aus gemessene] Länge eines Stücks von einer solchen kürzesten Linie sein und deren Endpunkte die Coordinaten $x,y,.z$ haben. Da $\theta$ und $s$ also einem ganz bestimmten Punkt der krummen Fläche angehören, so kann man $x,y,z$ wie Functionen von $\theta$ und $s$ betrachten. Die Richtung des Elements von $s$ entspreche auf der Kugelfläche dem Punkte $\lambda$, dessen Coordinaten $\xi,\eta,\zeta$ sind, so dass man
\[
\xi=\frac{\partial x}{\partial s},\qquad\eta=\frac{\partial y}{\partial s},\qquad\zeta=\frac{\partial z}{\partial s}
\]
haben wird.

Die Endpunkte aller kürzesten Linien von gleicher Länge $s$ entsprechen einer krummen Linie, deren Länge $t$ heissen mag. Man kann offenbar $t$ als eine Function von $s$ und $\theta$ betrachten, und wenn die Richtung des Elements von $t$ auf der Kugelfläche dem Punkte $\lambda'$ entspricht, dessen Coordinaten $\xi',\eta',\zeta'$ sind, wird man
\[
\xi'.\frac{\partial t}{\partial\theta}=\frac{\partial x}{\partial\theta},\qquad\eta'.\frac{\partial t}{\partial\theta}=\frac{\partial y}{\partial\theta},\qquad\zeta'.\frac{\partial t}{\partial\theta}=\frac{\partial z}{\partial\theta}
\]
haben. Es ist folglich
